\documentclass[11pt,letterpaper]{article}
\usepackage[T1]{fontenc}
\usepackage{lmodern,amsmath,amssymb,amsthm,mathtools,microtype}
\usepackage{booktabs,array,enumitem,listings,xcolor,url,textcomp}
\usepackage[letterpaper,margin=1in]{geometry}
\usepackage[hidelinks]{hyperref}
\hypersetup{pdftitle={Improved Upper Bounds on the Size of Condorcet Domains},
 pdfauthor={Pavlos Moros and Soren Riis}}
\setlength{\emergencystretch}{3em}
\setlength{\parindent}{0pt}
\setlength{\parskip}{0.35em}
\setlist{itemsep=0.15em,topsep=0.3em,leftmargin=*}
\newtheorem{theorem}{Theorem}[section]
\newtheorem{lemma}[theorem]{Lemma}
\newtheorem{proposition}[theorem]{Proposition}
\newtheorem{corollary}[theorem]{Corollary}
\theoremstyle{definition}
\newtheorem{definition}[theorem]{Definition}
\newtheorem{example}[theorem]{Example}
\newcommand{\F}{\mathcal F}
\newcommand{\Sig}{\Sigma}
\newcommand{\maj}{\operatorname{maj}}
\newcommand{\fr}{\operatorname{fr}}
\newcommand{\fz}{\operatorname{fz}}
\newcommand{\lk}{\operatorname{lk}}
\newcommand{\N}[2]{#1\mathrm N_{#2}}
\lstset{language=Python,basicstyle=\fontsize{8}{9.6}\selectfont\ttfamily,
 keywordstyle=\bfseries,commentstyle=\itshape,stringstyle=\ttfamily,
 columns=fullflexible,keepspaces=true,showstringspaces=false,upquote=true,
 breaklines=true,breakatwhitespace=false,tabsize=4,
 aboveskip=0.5em,belowskip=0.5em}
\begin{document}
\begin{center}
{\Large\bfseries Improved Upper Bounds on the Size\\[2pt]
of Condorcet Domains}\par
\medskip
Pavlos Moros \qquad S\o ren Riis\par
\small Journal review draft, 1 November 2026
\end{center}
\begin{abstract}
We prove \(f(n)\leq32n f(n-1)/(11n+8)\) for the maximum size of a
Condorcet domain, giving exponential base \(32/11<3\).
An exact deletion count, a decomposition into adjacent-swap components
and a local counting inequality give the bound; the same local
inequality gives base \(8/3\) for peak-pit domains.
Small examples illustrate the structural counts. The mathematical
argument is presented here; the exact finite calculations, reproducible
programs and complete Lean formalisation are supplied as accompanying
computational material.
\end{abstract}

\section{Introduction}\label{sec:introduction}
A Condorcet domain admits rankings for which pairwise majority is
transitive for every odd number of voters, allowing repeated rankings.
How many rankings can be admitted while retaining consistent majority
decisions? Write \(f(n)\) for the maximum on \(n\) alternatives.
We use Sen's equivalent combinatorial description: each triple has
an alternative forbidden from one of its three positions in every
admitted ranking~\cite{Sen}. A domain is \emph{peak-pit} when these
forbidden positions can always be chosen first or last. The alternative
and the position may vary with the triple; no common axis is required.
Write \(h(n)\) for the peak-pit maximum.

Fishburn~\cite{Fishburn} surveys the extremal problem, and
Raz~\cite{Raz} proved an exponential upper bound for \(f(n)\).
For peak-pit constructions and recent lower bounds, see Karpov and
Slinko~\cite{KarpovSlinko} and Karpov, Markstr\"om, Riis and
Zhou~\cite{KMRZ}. An earlier AI-assisted working draft of this
project~\cite{Moros} recorded the exact fibre count, the gap inequality,
and bounds with general bases \(56/19\) and \(1336/455\), and peak-pit
base \(8/3\). Its arguments and verification programs were developed
with Claude and examined with Codex; its finite correction coefficients
were selected by linear programming. The present paper consolidates
that investigation, adds the component reduction, and improves the
general base to \(32/11=2.90909\ldots\).

\begin{theorem}[Main bounds]\label{thm:main}
For every integer \(n\geq3\),
\[
 f(n)\leq\frac{32n}{11n+8}f(n-1),\qquad
 h(n)\leq\frac83h(n-1).
\]
Consequently \(f(n)\leq4(32/11)^{n-3}\),
\(h(n)\leq4(8/3)^{n-3}\), and
\(f(n)=O(n^{-8/11}(32/11)^n)\).
The limits \(\lim_{n\to\infty}f(n)^{1/n}\) and
\(\lim_{n\to\infty}h(n)^{1/n}\) exist and are at most \(32/11\)
and \(8/3\), respectively.
\end{theorem}

After completion, partition the domain into components under swaps
of consecutive entries. Each component is described by never-first/never-last
restrictions and fixed comparisons; its deletion count treats each
consecutive range of insertion positions separately. A majority projection,
related to the median structures of Puppe and Slinko~\cite{PuppeSlinko},
proves this directly. Combining the component counts with the gap
inequality gives the general bound. The only computational ingredient is
a local inequality. The complete Lean formalisation, finite data,
programs and verification instructions are provided in the accompanying
repository~\cite{LeanRepository}.

\clearpage
\section{Orders, triple restrictions and full families}\label{sec:objects}
Let \(X\) be a finite set of \(n\) alternatives. An \emph{order} is
a strict total order on \(X\), written from first to last; \(a>b\)
means that \(a\) occurs before \(b\). Restricting
an order to a subset means erasing all other labels; this restriction is
also called its \emph{trace} on the subset.
For a triple \(T\subseteq X\) and \(x\in T\), the \emph{never-condition}
\(\N{x}{i}\), for \(i\in\{1,2,3\}\), says that \(x\) does not occupy position \(i\) in the
restriction to \(T\). Thus \(\N{x}{1}\), \(\N{x}{2}\) and
\(\N{x}{3}\) mean never first, never middle and never last,
respectively, \emph{within that triple}. For instance, \(\N23\) on
\(\{1,2,3\}\) permits exactly \(123,213,231,321\).
A condition is \emph{common} to a family if every order satisfies it.
We also call a never-condition a \emph{law}.

\begin{definition}[Condorcet and peak-pit domains]\label{def:cd}
A \emph{Condorcet domain} is a family of orders with a common
never-condition on every triple. It is \emph{peak-pit} if each triple
has a common never-first or never-last condition. We call these two
kinds of condition \emph{extreme}; never-middle is the remaining kind.
\end{definition}
By Sen's characterisation~\cite{Sen}, this is equivalent to the voting
description in the introduction. Appendix~\ref{app:voting} includes
an elementary proof. The functions \(f(n),h(n)\) are attained maxima
because the set of orders is finite. A Condorcet domain is \emph{maximal}
if no further order can be added while preserving the Condorcet property;
a \emph{maximum} domain has largest possible cardinality. The first term concerns adding orders; the second concerns
the largest size among all domains on the given alternatives.

\begin{definition}[Full families and \(\F_n\)]\label{def:full}
Choose one never-condition for each triple, and collect \emph{every}
order satisfying all the chosen conditions. We call the resulting
family a \emph{full domain} for that choice of conditions.
To define \(\F_n\), use the labels \(0,\ldots,n-1\), choose an
extreme condition on every triple, and optionally require specified
pair comparisons, such as \(a>b\). The nonempty families consisting
of \emph{all} orders obeying such a system form \(\F_n\).
\end{definition}
Here ``full'' refers to satisfying a chosen system exactly; it does not
mean maximal. It also differs from the literature's ``closed'', meaning
closed under majority~\cite{PuppeSlinko}. Families on other label sets
are tested for membership in \(\F_n\) after relabelling by a bijection;
all conditions and counts are invariant under this change of names.

Every maximal Condorcet domain is full: choose a common condition on
each triple. All orders obeying the choices form a Condorcet enlargement,
so maximality forces equality. A maximum peak-pit domain likewise equals
its extreme-condition completion and belongs to \(\F_n\).
This completion need not be maximal among all Condorcet domains.
Whenever we delete a label we assume \(n\geq1\).

For any family \(D\), the \emph{swap graph} has its orders as vertices;
two are joined when exchanging consecutive entries takes one to the
other. Such a swap is \emph{enabled}. A path is a sequence of these
swaps staying in \(D\). A \emph{swap component} consists of all orders
reachable from a given order, including itself, by such paths.

\clearpage
\section{Insertions and the exact count}\label{sec:fibres}
For a family \(D\) on \(X\), with \(|X|=n\), define
\[
 D_{-x}=\{\sigma\text{ with }x\text{ erased}:\sigma\in D\},
 \qquad\Sig(D)=\sum_{x\in X}|D_{-x}|.
\]
Repeated outputs are counted only once in each \(D_{-x}\).
If \(D\) is a Condorcet domain, restriction preserves its triple laws, so
\begin{equation}\label{eq:deleteupper}
 \Sig(D)\leq nf(n-1),\qquad
 \Sig(D)\leq nh(n-1)\quad\hbox{if \(D\) is peak-pit}.
\end{equation}
For \(P\in D_{-x}\), write \(P[x:t]\) for inserting \(x\) with
\(t\) labels above it; \(t\) indexes its insertion slot. Its \emph{insertion fibre} is
\[
 F_D(x,P)=\{t\in\{0,\ldots,n-1\}:P[x:t]\in D\}.
\]
A \emph{run} is a maximal consecutive block of occupied slots.
A fibre with \(r\) runs contributes \(r-1\) to the total gap count
\(G(D)\); a long empty stretch between runs is one gap.
Let \(A(D)\) count pairs \((\sigma,j)\), where \(\sigma\in D\) and
the swap at positions \(j,j+1\), \(1\leq j<n\), is enabled.
Thus each edge counts twice. For a fixed family \(D\),
we abbreviate these statistics to
\(M=|D|,\ \Sig=\Sig(D),\ G=G(D),\ A=A(D)\).

The starting point is an exact count, with no inequality or lost term:
every order/label choice is accounted for by a deletion output, a gap
or an enabled-swap incidence.
\begin{lemma}[Fibre identity]\label{lem:fibre}
For every family of orders,
\(\boxed{\Sig+G+A=nM}\).
\end{lemma}
\begin{proof}
A fibre with \(m\) occupied slots and \(r\) runs has \(m-r\)
neighbouring occupied pairs. Summing \(m\) gives \(nM\), and summing
\(r\) gives \(\Sig+G\). Each neighbouring occupied pair is a swap
edge with the inserted label distinguished. Every edge occurs twice,
once for each exchanged label, giving exactly \(A\).
\end{proof}

\begin{example}[Deletion collisions]\label{ex:four}
The full domain \(D=\{123,213,231,321\}\) satisfies \(\N23\).
Its deletion outputs are
\[
 D_{-1}=\{23,32\},\qquad D_{-2}=\{13,31\},\qquad
 D_{-3}=\{12,21\}.
\]
For instance, deleting 1 merges \(123,213,231\) into 23, and
\(F_D(1,23)=\{0,1,2\}\). The swap graph is the path
\(123\leftrightarrow213\leftrightarrow231\leftrightarrow321\).
Thus \(M=4,\Sig=6,G=0,A=6\), and \(6+0+6=3\cdot4\).
\end{example}

The \emph{top slice at \(x\)} retains only orders starting with \(x\)
and then erases \(x\). The \emph{bottom slice at \(x\)} does the same
for orders ending with \(x\). A slice can be empty. A sum over slices
runs over the choices of fixed endpoint label.
\begin{lemma}[Two closure properties]\label{lem:interval}
Families in \(\F_n\) have interval insertion fibres (one run). Every nonempty
slice obtained by fixing and removing a first or last label belongs
to \(\F_{n-1}\).
\end{lemma}
\begin{proof}
Fix the relative order of the other labels. An extreme law permits
an interval of slots for the inserted label: never-first puts it
below the higher of the two others, and never-last above the lower.
If another label is distinguished, the law is automatic or puts the
inserted label on one specified side of it. Fixed comparisons also
give intervals. Intersect these intervals and use fullness.
For an endpoint slice, internal triple laws remain. A triple containing
the removed endpoint reduces to a fixed comparison, a tautology or a
contradiction. Pair constraints reduce similarly; contradictions mean
the slice is empty. The remaining constraints describe the slice exactly.
\end{proof}

\clearpage
\section{Reducing a general domain to extreme laws}\label{sec:components}
Throughout this section \(D\) is full. Its swap components are the
connected pieces of the graph defined in Section~\ref{sec:objects}. For
three orders write \(\maj(\alpha,\beta,\gamma)\) for their pairwise
majority order. The majority-closure argument is part of the median
structure studied by Puppe and Slinko~\cite{PuppeSlinko}; we give the
specific adjacent-swap component argument in full.

\begin{lemma}[Majority projection]\label{lem:projection}
The majority of three members of \(D\) lies in \(D\).
If \(K\) is a swap component of \(D\) and \(U\) is the set of all oriented
pair comparisons common to every order of \(K\), then
\[
 K=\{\tau\in D:\tau\text{ satisfies every comparison in }U\}.
\]
\end{lemma}
\begin{proof}
Majority is an order by Appendix~\ref{app:voting}. On a triple
\(\{x,y,z\}\), the law \(\N x2\) says \(x>y\) if and only if
\(x>z\); the two majority comparisons therefore agree as well.
If majority violated a never-first law, two sets of at least two
voters among three would support its two comparisons. Their intersection
would contain a voter violating the law. Never-last is symmetric.
Fullness therefore puts the majority order in \(D\).

For the displayed equality, suppose \(\tau\in D\) satisfies \(U\).
Choose \(\sigma\in K\) disagreeing with \(\tau\) on as few pair
comparisons as possible. If they differ, a differing comparison is not
fixed on \(K\). Take \(\rho\in K\) agreeing with \(\tau\) there,
and project a path from \(\sigma\) to \(\rho\) by
\(v\mapsto\maj(\sigma,\tau,v)\). Images lie in \(D\); consecutive
images differ in at most one comparison. Two distinct orders differing
in one comparison exchange adjacent labels: an intervening label would
force another comparison to change. Remove consecutive repeated images;
the remaining path stays in \(K\) and starts at \(\sigma\). On comparisons where \(\sigma,\tau\)
agree it is constant; elsewhere it follows \(v\). Its last order
corrects the selected disagreement without losing any agreement.
This contradicts the choice of \(\sigma\), proving the claim.
\end{proof}

\begin{proposition}[Component identity]\label{prop:components}
Every component \(K\) belongs to \(\F_n\), and
\begin{equation}\label{eq:components}
 \sum_K |K|=M,\qquad \sum_K\Sig(K)=\Sig(D)+G(D).
\end{equation}
\end{proposition}
\begin{proof}
A chosen never-middle law \(\N x2\) equates the two comparisons of
\(x\) on that triple. Neither can change along a swap path, which
changes only one comparison at a time. Thus throughout \(K\), \(x\)
is above both others or below both. Replace the law by \(\N x3\)
or \(\N x1\), respectively. Keep the original extreme laws and add
all comparisons \(U\) fixed on \(K\). Every order in \(K\) satisfies
these constraints. Conversely, they recover the old never-middle laws,
so an order satisfying them lies in \(D\) and agrees with \(U\).
Lemma~\ref{lem:projection} puts it in \(K\). Thus these conditions
describe \(K\) exactly, proving \(K\in\F_n\).

Components partition \(D\), giving the first identity. Each original
fibre run lies in one component, because consecutive occupied slots
are swap edges. A component cannot contain two separated runs: its
own fibre is an interval by Lemma~\ref{lem:interval}. Thus each original
run gives exactly one component deletion output. Summing runs gives
the second identity.
\end{proof}

\begin{example}[Counting separated runs]\label{ex:gap}
The full domain \(D=\{123,132,231,321\}\), defined by \(\N12\),
has components \(\{123,132\}\) and \(\{231,321\}\).
Its 1-fibres above 23 and 32 both have slots \(\{0,2\}\), each
contributing one gap. Thus \(\Sig(D)=6,G(D)=2\). Each component
has deletion sum 4, so \(4+4=6+2\): splitting the family counts
its separated insertion runs separately.
\end{example}

\clearpage
\section{The gap inequality}\label{sec:gap}
Fix a Condorcet domain \(D\) on \(X\), with \(|X|=n\).
A \emph{lift} moves one label upwards, and a \emph{drop} moves it
downwards, keeping all other relative positions and remaining in \(D\).
These moves may jump several slots. An incidence \((\sigma,x)\) is
\emph{free} if it has both moves, and \emph{frozen} if it has neither;
their counts are \(\fr(D),\fz(D)\). In an order \(\sigma\), consecutive
entries \(a,c\) form a \emph{locked} pair if \(a\) has no drop and
\(c\) has no lift. Let \(\lk(D)\) count these triples \( (\sigma,a,c)\). With \(D\) fixed, abbreviate these
counts to \(\fr,\fz,\lk\).

\begin{lemma}[Two counts]\label{lem:counts}
For a full domain,
\begin{equation}\label{eq:movecounts}
 2\Sig=nM-\fr+\fz=(n+1)M-A+\lk,
 \qquad 3\Sig+\fr=(n+1)M+G+\fz+\lk.
\end{equation}
\end{lemma}
\begin{proof}
Each fibre has one order with no lift and one with no drop. Their
combined count is \(2\Sig\); an incidence contributes 0, 1 or 2
according as it is free, has one move, or is frozen.
For an adjacent pair, the missing drop of its upper label and missing
lift of its lower label contribute 0 if the swap is enabled, 2 if
locked, and 1 otherwise. Indeed a blocked swap changes only that pair
comparison and violates a chosen law on a triple containing the pair.
Write its old trace as \(zac\) if the third label is above. The
violated law is \(\N a3\) or \(\N c2\), either of which forbids
every drop of \(a\). If the third label is below, the old trace is
\(acz\); the law is \(\N a2\) or \(\N c1\), forbidding every
lift of \(c\). Thus both crossing moves cannot be available.
The first label also has no lift and the last no drop. The second
count is therefore \(2M+(n-1)M-A+\lk\). Combine the two counts
with \(\Sig+G+A=nM\) for the final identity.
\end{proof}

\begin{lemma}[Free-to-locked injection]\label{lem:injection}
For a maximal Condorcet domain, \(\fr\leq\lk\).
\end{lemma}
\begin{proof}
First, if \(a>c\) in every order, they cannot be adjacent in one.
Swapping such an occurrence adds an order while preserving a common
law on every triple: triples containing the pair retain \(\N c1\)
if the third label is above, or \(\N a3\) if below. All other triples
are unchanged. This contradicts maximality.

A free label \(b\) is neither first nor last. Write \(a,b,c\) for the
consecutive block centred on \(b\). The order, a lift
and a drop give \(abc,bac,acb\); their only possible common laws
are \(\N a3\) and \(\N c1\). Both cannot hold, since together
they force \(a>c\) everywhere, whereas a lift puts them consecutively.
If \(\N a3\) holds, take the available lift leaving \(b\) as low
as possible; otherwise take the available drop leaving it as high
as possible: use the closest available slot in the chosen direction.
In the image order the pair
\(a,c\) is consecutive and locked: either crossing move would violate
the selected law. Map the free incidence to this order/pair incidence.
For a specified moved label, the preimage is recovered uniquely by
putting it back between the image pair. Appendix~\ref{app:collisions}
excludes collisions between different moved labels, completing the proof.
\end{proof}

The first identity in Lemma~\ref{lem:counts} is the sum over deleted
labels of Fathi's fibre-excess identity~\cite[eq.\ (96)]{Fathi}.
The two possible laws at a free incidence are also Fathi's flank-rule
observation~\cite[Lemma K.52]{Fathi}. The injection uses maximality
to make the assignment collision-free.

\begin{corollary}[Gap inequality]\label{cor:gap}
For every maximal Condorcet domain,
\(G+(n+1)M\leq3\Sig\).
\end{corollary}
\begin{proof}
In the last identity of~\eqref{eq:movecounts}, use
\(\fr\leq\lk\) and \(\fz\geq0\).
\end{proof}

\clearpage
\section{A precise local inequality}\label{sec:finite}
The next ingredient concerns seven consecutive positions. Its
correction terms depend only on six-label endpoint slices and cancel
when windows overlap. Exact data and exhaustive checks are supplied in
the computational material~\cite{LeanRepository}.
Recall that top and bottom slices fix and delete the first and last
label, respectively. Empty slices contribute zero.
For \(J\in\F_7\), let \(a_j(J)\) count orders whose swap at
positions \(j,j+1\) is enabled, for \(1\leq j\leq6\). Define
\begin{equation}\label{eq:cost}
 C(J)=\frac{a_1(J)+2a_2(J)+2a_3(J)+2a_4(J)+2a_5(J)+a_6(J)}{10}.
\end{equation}
Summing \(A(H)\) over top and bottom six-slices counts the end swaps
once and all other swaps twice; hence that sum is \(10C(J)\).

\begin{lemma}[Seven-position inequality]\label{lem:finite}
There is a fixed rational-valued function \(q\), defined below, such that
every \(J\in\F_7\), and the empty family, satisfies
\begin{equation}\label{eq:finite}
 C(J)\leq\frac58|J|+
   \sum_{\text{bottom }H}q(H)-\sum_{\text{top }H}q(H).
\end{equation}
\end{lemma}

On five labels, enumerate the families in \(\F_5\),
identify families related by a bijective renaming of labels
(\emph{isomorphic} families), and call each such class a \emph{type}. Encode a family by a binary \emph{mask}:
list all permutations in lexicographic order, starting at index 0,
and set bit \(i\) exactly when permutation \(i\) belongs to the family.
Order the types by their least masks over all label renamings.
There are 1160 types,
indexed \(0,\ldots,1159\). Set \(\phi(P)=W_k/d\) on type \(k\), where
\(d=8\,846\,117\,280\). The complete coefficient table and its
exhaustive generation are supplied with the computation~\cite{LeanRepository}.

For a six-label family \(H\in\F_6\), define
\begin{align*}
 b(H)&=\sum_{\text{bottom }P}|P|\phi(P),&
 t(H)&=\sum_{\text{top }P}|P|\phi(P),\\
 s(H)&=\frac{12}{19}|H|+b(H)-t(H)-\frac15A(H),&
 \eta&=\frac1{152}.
\end{align*}
Let \(r_b(H)\) count bottom five-slices of type 866, and
\(r_t(H)\) top five-slices of type 41. The accompanying table
specifies fourteen exceptional six-types: five assigned to \(\mathcal T\)
and nine to \(\mathcal S\). Both classes include every relabelled copy.
The nine integer corrections define \(dp(H)\) on \(\mathcal S\).
Outside \(\mathcal S\), put \(p(H)=2\eta(r_t(H)-r_b(H))\). Finally set
\[
 \theta_b(H)=s(H)/2+p(H),\qquad\theta_t(H)=s(H)/2-p(H),
 \qquad q(H)=(b(H)+t(H))/2+p(H).
\]
All these quantities are zero on the empty family. Arguments are
suppressed when a single slice \(H\) is under discussion.

The following facts are verified by exhaustive exact computation;
the coverage proof and code are supplied separately~\cite{LeanRepository}.
\begin{enumerate}[label=F\arabic*.]
\item There is a unique eight-order family \(R_8\in\mathcal T\)
      up to relabelling. All its copies \(H\) have
      \(\theta_b(H)=\theta_t(H)=0\).
\item Outside \(\mathcal T\cup\mathcal S\),
 \(\theta_b\geq\eta|H|/2+4\eta r_t\) and
 \(\theta_t\geq\eta|H|/2+4\eta r_b\).
\item If \(J\) has a six-slice in \(\mathcal S\cup\mathcal T\)
 that is not isomorphic to \(R_8\), then
\[
 \Lambda(J):=\sum_{\text{bottom }H}\theta_b(H)+
             \sum_{\text{top }H}\theta_t(H)\geq\eta|J|.
\]
\end{enumerate}

\clearpage
\section{Paying for the remaining slices}\label{sec:receiver}
We now prove the local inequality from the explicitly checked finite
facts; the remaining argument is by hand. A convenient labelled
representative of the eight-order exceptional type is
\[
 R_8=\{524130,524310,534120,534210,
       542130,542310,543120,543210\}.
\]
Its common first and last labels are 5 and 0. Deleting the first
gives five-type 866, and deleting the last gives type 41.
Its only nonempty bottom and top slices therefore have types 41 and
866, respectively, so \(r_b(R_8)=r_t(R_8)=0\).
These statements and its equality with the supplied eight-order type
are also verified by the indexing procedure supplied with the computation.

\begin{proof}[Proof of Lemma~\ref{lem:finite}]
F3 already gives \(\Lambda(J)\geq\eta|J|\) unless all nonempty
six-slices are regular (outside \(\mathcal T\cup\mathcal S\)) or
of type \(R_8\). Consider this remaining case.

Each regular slice contributes a baseline \(\eta|H|/2\) by F2.
An \(R_8\) slice contributes zero by F1 and is missing baseline
\(4\eta\). Suppose a bottom \(R_8\) slice is obtained by fixing
last label \(z\). Its orders share first label \(w\), and therefore
also enter the top slice at \(w\). This top receiver, upon fixing
and deleting its last label \(z\), has a five-slice of type 866.
Thus its \(r_b\) is positive. It cannot be \(R_8\), whose \(r_b\)
is zero, so it is regular. F2 supplies an extra \(4\eta\) in its
\(\theta_t\), exactly the missing baseline.

Different bottom \(R_8\) slices have different last labels. They
charge different receiver/last-label pairs, each counted separately
in \(r_b\), so no contribution is reused. Reverse top and bottom
for top \(R_8\) slices; their receivers use the separate
\(\theta_b\) contributions from \(r_t\). All missing baselines
are paid. At each end the slice sizes sum to \(|J|\), yielding
total baseline \(\eta|J|\). Thus the bound on \(\Lambda\) holds
in every case.

It remains to expand the definitions. Since the slice sizes sum to
\(|J|\) at each end,
\begin{align*}
 \Lambda(J)
 &=\frac{12}{19}|J|-\frac1{10}
       \sum_{\text{top and bottom }H}A(H)\\
 &\quad+\frac12\sum_{\text{bottom and top }H}(b(H)-t(H))
       +\sum_{\rm bottom}p(H)-\sum_{\rm top}p(H).
\end{align*}
The sum of \(t(H)\) over bottom slices equals the sum of \(b(H)\)
over top slices. Both enumerate the same conditional five-families
obtained by fixing one first and one last label and deleting them.
The deletions commute, including their label normalisations. Hence
\[
 \Lambda(J)=\frac{12}{19}|J|-C(J)
       +\sum_{\rm bottom}q(H)-\sum_{\rm top}q(H).
\]
Use \(\Lambda(J)\geq\eta|J|\) and
\(12/19-1/152=5/8\). For the empty family every term is zero.
\end{proof}

The exceptional case is therefore resolved by a counting argument,
rather than by checking all seven-label families. The computation
handles only the thirteen other exceptional six-types and their two
endpoint orientations. The accompanying computation gives a full coverage proof.

\clearpage
\section{From local windows to an exact density}\label{sec:windows}
\begin{proposition}\label{prop:density}
For every \(n\geq1\) and \(K\in\F_n\),
\(8A(K)\leq5n|K|\) and \(\Sig(K)\geq3n|K|/8\).
\end{proposition}
\begin{proof}
First suppose \(n\geq7\), and write \(M=|K|\).
At each block of seven consecutive positions, fix the complete ordered
prefix and suffix outside it. Its conditional family \(J\) belongs
to \(\F_7\) when nonempty, by repeated endpoint slicing. At each
window position these families partition \(K\), so their sizes sum
to \(M\). A swap within a window has the same availability in its
conditional family as in \(K\). Apply~\eqref{eq:finite} to all contexts
and windows.

The top six-slices of one window are exactly the bottom six-slices
of the next: each fixes the same ordered exterior and retains the
same central six labels. For example the windows \(xabcdef\) and
\(abcdefy\) share \(abcdef\), with \(x\) and \(y\) fixed outside
after slicing. Thus all internal \(q\)-terms cancel. This matching
retains the full exterior, not merely the isomorphism type of a slice.

The class \(\F_6\) is finite, so
\(Q=\max(\{0\}\cup\{|q(H)|/|H|:H\in\F_6\})\) is finite.
Thus \(|q(H)|\leq Q|H|\), also for the empty family.
Each of the two surviving outside layers has total size
\(M\), so its absolute contribution is at most \(QM\). Consequently
\[
 \sum_{\text{all window contexts }J}C(J)
 \leq\frac58(n-6)M+2QM.
\]
A swap contained in six windows has total weight
\((1+2+2+2+2+1)/10=1\). Every swap has weight between zero and one.
The sum of all position weights is \(n-6\), whereas there are \(n-1\)
swap positions. The total deficit is exactly five per order. Since
each swap indicator is at most one,
\begin{equation}\label{eq:affine}
 A(K)\leq\frac58nM+BM,\qquad B=5/4+2Q.
\end{equation}

Fix any \(K\in\F_n\), now allowing small \(n\), and concatenate
\(t\) disjoint copies in a fixed block order. Each block chooses an
order independently; all comparisons between blocks are fixed.
A mixed triple has a label fixed first or last, so compatible extreme
laws can be chosen. The resulting full family \(K^{[t]}\) belongs
to \(\F_{tn}\), and
\[
 |K^{[t]}|=M^t,\qquad A(K^{[t]})=tA(K)M^{t-1}.
\]
The second identity counts an enabled swap in any one block and
independent choices in the others. A swap across a join reverses a
fixed comparison and is unavailable; it creates no new term.
For \(tn\geq7\), apply~\eqref{eq:affine} to the joined family:
\[
 tA(K)M^{t-1}\leq\frac58tnM^t+BM^t,
 \qquad A(K)\leq\left(\frac58n+\frac Bt\right)M.
\]
Let \(t\) grow. This proves \(A(K)\leq5nM/8\), including the
original small \(n\). Its fibres are intervals, so \(G=0\);
Lemma~\ref{lem:fibre} then gives the deletion bound.
\end{proof}

\begin{example}[Two joined copies]\label{ex:copies}
Join the four-order domain \(D\) of Example~\ref{ex:four} to its
relabelled copy \(D'=\{456,546,564,654\}\), always putting its first
block above its second, and call the joined family \(J\in\F_6\).
It has 16 orders. Each deletion leaves two choices in its affected
block and four in the other, so \(\Sig(J)=6\cdot2\cdot4=48\).
The internal paths give \(A(J)=2\cdot6\cdot4=48\), and \(G(J)=0\).
The join adds no swap: crossing it reverses a fixed comparison.
\end{example}

\clearpage
\section{The upper bounds and their scope}\label{sec:finish}
\begin{proof}[Proof of Theorem~\ref{thm:main}]
Let \(D\) be a maximum Condorcet domain on \(n\) labels, with
\(M=f(n)\). It is maximal and therefore full. By the component
identity and Proposition~\ref{prop:density},
\[
 \frac38 nM\leq\sum_K\Sig(K)=\Sig+G.
\]
The gap inequality gives \(\Sig+G\leq4\Sig-(n+1)M\). Combining,
\[
 \frac38nM\leq4\Sig-(n+1)M,
 \qquad (11n+8)M\leq32\Sig\leq32n f(n-1).
\]
This proves the general recurrence. A maximum peak-pit domain belongs
to \(\F_n\), so its deletion sum is at least \(3nM/8\) and at
most \(nh(n-1)\), proving the peak-pit recurrence.

On three labels a common law omits two of the six orders; the
four-order domain in Example~\ref{ex:four} attains four. Thus
\(f(3)=h(3)=4\). Iterating the recurrences gives both geometric
bounds, since \(32n/(11n+8)\leq32/11\).
More precisely, with \(r=32/11\) and \(a=8/11\),
\[
 f(n)\leq4r^{n-3}\prod_{k=4}^{n}\frac{k}{k+a}
 \leq4r^{n-3}\left(\frac4{n+1}\right)^a.
\]
For the last step, \((1+u)^a\leq1+au\) for \(u\geq0\):
\(1+au-(1+u)^a\) is zero at zero and has nonnegative derivative.
Thus \(k/(k+a)\leq(k/(k+1))^a\), and the product telescopes.
This proves the claimed polynomial correction.

Finally, concatenate two domains in a fixed block order. Majority is
transitive within each block and respects the block order; mixed
triples have a fixed extreme label. Thus both the Condorcet and
peak-pit properties are preserved, and cardinalities multiply.
Hence for \(g=f\) or \(h\),
\(g(m+n)\geq g(m)g(n)\). Singletons give \(g(n)\geq1\).
Put \(a_n=\log g(n)\geq0\); the unique empty order gives
\(g(0)=1\). For fixed \(k\geq1\), write \(n=qk+s\),
\(0\leq s<k\). Then \(a_n\geq q a_k+a_s\geq q a_k\), so
\(\liminf a_n/n\geq a_k/k\). Taking the supremum over \(k\)
shows that the liminf is at least \(\sup_k a_k/k\); the limsup
cannot exceed it. The geometric bound makes this supremum finite
and bounds the limit by \(\log r\), or \(\log(8/3)\).
Exponentiating proves the root-growth assertions.
\end{proof}

\subsection*{Computation and formal verification}
The structural counts, component decomposition, receiver argument and
limiting steps are hand proofs; Appendix~\ref{app:collisions} supplies
all collision cases of the injection. The sole computational input is
the finite comparison in Section~\ref{sec:finite}. Its coefficients,
exhaustive coverage argument, exact integer checks and reproduction
instructions are supplied with the paper and in the repository~\cite{LeanRepository}.
All mathematical comparisons use exact integers and no optimiser.

All seven conclusions of Theorem~\ref{thm:main} are also proved end to
end in Lean~4.26.0 with Mathlib~4.26.0. The entry point is
\texttt{CDCert.Main}; the local inequality is \texttt{CDCert.K6.finite\_lemma}.
Finite certificates and their coverage are checked by Lean's kernel:
a successful Python run is not a premise of the formal proof.
The final results use no unproved placeholders or additional axioms,
only Lean's standard \texttt{propext}, \texttt{Classical.choice} and
\texttt{Quot.sound}. The repository records the exact build, theorem
mapping and verification scope~\cite{LeanRepository}; this verifies the
mathematical statements, not historical claims or the Python programs themselves.

Claude and ChatGPT/Codex assisted this work; the mathematical claims
and their presentation remain the authors' responsibility.

\clearpage
\begin{thebibliography}{9}
\small\raggedright
\setlength{\itemsep}{4pt}
\bibitem{Sen} A.~K.~Sen, A possibility theorem on majority decisions,
\emph{Econometrica} \textbf{34} (1966), 491--499.
\url{https://doi.org/10.2307/1909947}.
\bibitem{Fishburn} P.~C.~Fishburn, Acyclic sets of linear orders:
a progress report, \emph{Social Choice and Welfare} \textbf{19}
(2002), 431--447. \url{https://doi.org/10.1007/s003550100120}.
\bibitem{Raz} R.~Raz, VC-dimension of sets of permutations,
\emph{Combinatorica} \textbf{20} (2000), 241--255.
\url{https://doi.org/10.1007/s004930070023}.
\bibitem{KarpovSlinko} A.~Karpov and A.~Slinko, Constructing large
peak-pit Condorcet domains, \emph{Theory and Decision} \textbf{94}
(2023), 97--120. \url{https://doi.org/10.1007/s11238-022-09878-9}.
\bibitem{KMRZ} A.~Karpov, K.~Markstr\"om, S.~Riis and B.~Zhou,
\emph{Improved lower bounds for the maximum size of Condorcet domains},
preprint (2026). \url{https://arxiv.org/abs/2601.07336}.
\bibitem{Moros} P.~Moros, \emph{Condorcet domains have fewer than
\(2.95^n\) orders}, private manuscript, version \texttt{upper-v10},
23 October 2026. AI-assisted working draft from an earlier stage of
the present project.
\bibitem{PuppeSlinko} C.~Puppe and A.~Slinko, Condorcet domains,
median graphs and the single-crossing property,
\emph{Economic Theory} \textbf{67} (2019), 285--318.
\url{https://doi.org/10.1007/s00199-017-1084-6}.
\bibitem{Fathi} K.~Fathi, \emph{The peak-pit maximum for Condorcet
domains on nine alternatives: cube classifications and exact deletion
bounds \(h(9)=492\)}, preprint, Zenodo version v8
(manuscript Version 41, 5 September 2026), eq.\ (96) and Lemma K.52.
\url{https://doi.org/10.5281/zenodo.22346663}.
\bibitem{LeanRepository} P.~Moros and S.~Riis,
\emph{Improved Upper Bounds on the Size of Condorcet Domains:
Lean formalisation and computational verification}, accompanying repository,
review snapshot \texttt{0d5dd168854b} (29 October 2026).
\url{https://github.com/SR123/condorcet-upper-bound}.
The Lean source and build records are pinned to this snapshot;
updated expository and standalone computational material accompanies
this review draft. Archival-release DOI to be added after author agreement.
\end{thebibliography}

\clearpage
\appendix
\section{Voting interpretation and the injection}\label{app:elementary}
\subsection{Equivalence with the voting definition}\label{app:voting}
A voter profile is a finite list of admitted orders, with repetitions
allowed. Pairwise majority puts \(a\) above \(b\) if more than half
the voters do. We prove that Definition~\ref{def:cd} is equivalent to
transitivity of majority for every odd positive profile size.

On a triple, each of \(\{abc,bca,cab\}\) and
\(\{acb,cba,bac\}\) gives a majority cycle. For a voting Condorcet
domain, at least one order in each set must therefore be missing from
its triple restrictions. By cyclic relabelling let the first missing
order be \(abc\). Pairing it with \(acb,cba,bac\) gives, respectively,
\(\N a1,\N b2,\N c3\) on all the remaining orders.
Conversely, under \(\N x2\) the two comparisons of \(x\) agree for
each voter, so they agree for majority, excluding a cycle.
Under \(\N x1\), each voter with \(x>y\) has \(z>x>y\).
A majority for \(x>y\) therefore forces a majority for \(z>y\),
excluding \(x>y>z>x\). Exchange \(y,z\) for the reverse cycle,
and reverse orders to treat \(\N x3\). An odd majority compares
every distinct pair in exactly one direction. If it is not transitive,
there are \(a>b\), \(b>c\), but \(c>a\): a directed triangle.
Thus excluding cycles on triples proves transitivity. This proves the
equivalence attributed to Sen~\cite{Sen}.

\subsection{The injection: all collision cases}\label{app:collisions}
We complete the hand proof of Lemma~\ref{lem:injection}.
Let the target locked pair be \(p,q\) in its image order \(\tau\).
For a preimage with moved label \(b\), call it \emph{upper} if \(b\)
lies above \(p\) in \(\tau\), and \emph{lower} otherwise.
For an upper preimage, \(D\) has common law \(\N p3\) on
\(\{p,q,b\}\), whereas \(\N q1\) is not common to \(D\). There is consequently an order \(\chi_b\) with
trace \(qpb\): the alternative \(qbp\) would violate \(\N p3\).
For a lower preimage, \(\N q1\) is common to \(D\), and \(\N p3\) is not.
The original neighbours in either preimage are \(p,b,q\).

\subsection*{Two upper preimages}
Suppose different moved labels \(b_1,b_2\) give the same target.
Let their preimage orders be \(\sigma_1,\sigma_2\), and drop
\(b_i\) in \(\sigma_i\) to obtain \(\rho_i\); such drops exist
because the preimages are free. Their traces on \(\{b_1,b_2,q\}\) are
\[
\begin{array}{c|cccc}
 \text{order}&\sigma_1&\sigma_2&\rho_1&\rho_2\\\hline
 \text{trace}&b_2b_1q&b_1b_2q&b_2qb_1&b_1qb_2.
\end{array}
\]
Both \(b_1,b_2\) occupy all three positions among these orders;
\(q\) occupies the last two. Thus the only common law is \(\N q1\).
In \(\chi_{b_1}\), \(q\) is above \(p,b_1\). The law just
found forces \(b_2\) above \(q\), giving \(b_2qp\) on
\(\{p,q,b_2\}\), contrary to its law \(\N p3\).

\subsection*{Two lower preimages}
Reverse every order. Majority transitivity and maximality are preserved.
Lifts and drops exchange roles, and the unique law used by the assignment
changes from never-first to never-last. The highest drop becomes the
lowest lift. Thus the reversed assignment has two upper preimages,
which the preceding argument excludes.

\subsection*{One upper and one lower preimage}
Let \(b_1\) be upper and \(b_2\) lower. On the four indicated labels,
\[
 \tau:b_1pqb_2,\qquad
 \sigma_1:pb_1qb_2,\qquad
 \sigma_2:b_1pb_2q.
\]
Take \(\chi=\chi_{b_1}\), whose trace on \(\{p,q,b_1\}\) is
\(qpb_1\). The lower preimage's law \(\N q1\) on
\(\{p,q,b_2\}\) forces \(b_2\) above \(q\), so \(\chi\)
has trace \(b_2qpb_1\). On \(\{b_1,b_2,q\}\), the orders
\(\tau,\sigma_2,\chi\) give
\[
 b_1qb_2,\qquad b_1b_2q,\qquad b_2qb_1.
\]
These allow exactly two common laws: \(\N{b_1}2\) and \(\N q1\).
Drop \(b_1\) from \(\sigma_1\) to obtain \(\rho\).
Now \(q\) lies above \(b_1\), and its comparison with \(b_2\)
is unchanged, so it lies above both. This fourth witness excludes
\(\N q1\), leaving \(\N{b_1}2\). It forces the trace of
\(\rho\) to be \(qb_2b_1\).

Lift \(b_2\) from \(\sigma_2\) to obtain \(\lambda\). Its
comparison with \(q\) is now upwards and the comparison of
\(b_1,q\) is unchanged. The same never-middle law forces its trace
on \(\{b_1,b_2,q\}\) to be \(b_1b_2q\).
On \(\{p,b_1,b_2\}\), the five orders
\(\tau,\sigma_1,\rho,\lambda,\chi\) therefore give
\[
 b_1pb_2,\quad pb_1b_2,\quad pb_2b_1,\quad
 b_1b_2p,\quad b_2pb_1.
\]
Each label occupies every position, leaving no common never-law.
This contradicts Definition~\ref{def:cd}. All collisions are
excluded, proving injectivity.\hfill\(\square\)

\end{document}
